\section{Introduction}
\label{sec:intro}

Evaluating a language model on a benchmark means answering every item, and the cost multiplies when many checkpoints, prompts or inference settings are compared. A natural saving is to run the new model on a small subset of items and extrapolate to the full benchmark. Because per-item results of earlier models are routinely available, recent work uses them to choose the subset and to fit the extrapolation: clustering items into ``anchor points''~\citep{vivek2024anchor}, combining them with item response theory (IRT)~\citep{mauricio2024tiny,kipnis2025metabench}, or learning assessors and policies from earlier evaluations~\citep{pacchiardi2024100,li2025active}. These works report that tens to hundreds of items suffice, usually compared with random subsets.

The practical question is narrower than ``can few items predict the score''. A random subset already gives an unbiased estimate with a known standard error, $\sqrt{p(1-p)/n}$ up to a finite-population factor. Source models can help only by reducing variance or by encoding which items are informative, and they can hurt when the target differs from the sources, because a fitted model of the sources is then biased. We therefore ask, for binary item-level scores and at equal item budget $n\ll N$:
\begin{itemize}\setlength\itemsep{0pt}
\item[\textbf{Q1}] Does choosing items and fitting predictors with source-model data beat a random subset, in error (MAE, RMSE) and in ranking (Spearman, Kendall)?
\item[\textbf{Q2}] How does the answer change with the budget $n$, and how many items are needed for a given error?
\item[\textbf{Q3}] Does it survive when the target models are stronger than every source model?
\item[\textbf{Q4}] Is the error due to bias or to variance, and can models closer than the typical error be ranked?
\end{itemize}

\paragraph{Contributions.}
(i) A leakage-free comparison protocol on 13 public matrices in which every method gets identical sources, targets, budgets and seeds, hyperparameters are fixed in advance or chosen on source models only, and uncertainty is a paired bootstrap over seeds (Section~\ref{sec:setup}). (ii) Re-implementations of random, stratified, calibrated, anchor-point, IRT, low-rank and supervised-selection pipelines, plus a stratified design with a design-based residual correction (P3; the correction is the standard difference estimator of survey sampling, so P3 is a tested design rather than a new estimator) whose design was frozen by cross-validation inside the source models before the final run (Section~\ref{sec:methods}). (iii) The evidence that source-informed methods win clearly under random splits, that model-based ones lose to random sampling under extrapolation because of bias, and that designs with a design-based residual correction, and stratified sampling with the plain mean, are the only ones whose MAE is never significantly worse than random in our runs (also with intervals over target models), with a gain that is real but small; ablations show that the correction, not stratification, carries the robustness (Section~\ref{sec:exp}). The scope of these statements, which includes the implementation differences from the original methods, is stated in Section~\ref{sec:disc}.
